% book10.tex --- Book X of Euclid's Elements: Incommensurable Magnitudes.
%
% All 115 propositions encoded.  Book X is the longest of the Elements
% and develops a classification of irrational magnitudes via repeated
% quadratic operations: it covers commensurability (X.1-X.18), medials
% (X.19-X.35), and the thirteen irrational lines including binomials and
% apotomes (X.36-X.115).  Proof sketches are condensed; the dependency
% structure is what we capture.
%
% Wording follows Heath (1908).

\section{Book X --- Incommensurable Magnitudes}
\label{sec:book-X}

\begin{claim}[Proposition X.1: Method of exhaustion lemma]
\label{prop:X.1}
Two unequal magnitudes being set out, if from the greater there be
subtracted a magnitude greater than its half, and from that which is
left a magnitude greater than its half, and if this process be
repeated continually, there will be left some magnitude which will be
less than the lesser magnitude set out.
\end{claim}
\begin{evidence}[Proof of X.1]
\label{ev:X.1}
By the Archimedean property (Definition V.4), some multiple of the
smaller magnitude exceeds the larger.  Iterated halving (or more)
brings the remainder below the smaller in a finite number of steps.
\dependson{X.1}{def:V.4}
\end{evidence}

\begin{claim}[Proposition X.2: Anthyphairesis detects incommensurability]
\label{prop:X.2}
If, when the lesser of two unequal magnitudes is continually
subtracted in turn from the greater, that which is left never
measures the one before it, the magnitudes will be incommensurable.
\end{claim}
\begin{evidence}[Proof of X.2]
\label{ev:X.2}
A common measure would persist through anthyphairesis (Common Notion
3); non-termination of the algorithm thus implies no common measure
exists.
\dependson{X.2}{X.1}
\dependson{X.2}{cn:3}
\dependson{X.2}{def:X.1}
\end{evidence}

\begin{claim}[Proposition X.3: Greatest common measure of commensurables]
\label{prop:X.3}
Given two commensurable magnitudes, to find their greatest common
measure.
\end{claim}
\begin{evidence}[Proof of X.3]
\label{ev:X.3}
Apply anthyphairesis (the Euclidean algorithm on magnitudes); by X.2
the algorithm terminates exactly when a common measure exists.
\dependson{X.3}{X.2}
\dependson{X.3}{def:X.1}
\end{evidence}

\begin{claim}[Proposition X.4: GCM of three commensurables]
\label{prop:X.4}
Given three commensurable magnitudes, to find their greatest common
measure.
\end{claim}
\begin{evidence}[Proof of X.4]
\label{ev:X.4}
Apply X.3 to the first two; then to the result with the third.
\dependson{X.4}{X.3}
\end{evidence}

\begin{claim}[Proposition X.5: Commensurables have a number ratio]
\label{prop:X.5}
Commensurable magnitudes have to one another the ratio which a number
has to a number.
\end{claim}
\begin{evidence}[Proof of X.5]
\label{ev:X.5}
If $d$ is a common measure of $a$, $b$, then $a = md$, $b = nd$, so
$a : b = m : n$ (a ratio of integers).
\dependson{X.5}{def:V.5}
\dependson{X.5}{def:X.1}
\end{evidence}

\begin{claim}[Proposition X.6: Number ratio implies commensurable]
\label{prop:X.6}
If two magnitudes have to one another the ratio which a number has
to a number, the magnitudes will be commensurable.
\end{claim}
\begin{evidence}[Proof of X.6]
\label{ev:X.6}
Converse of X.5: if $a : b = m : n$, dividing $a$ by $m$ produces a
common measure of $a$ and $b$.
\dependson{X.6}{X.5}
\end{evidence}

\begin{claim}[Proposition X.7: Incommensurables have no number ratio]
\label{prop:X.7}
Incommensurable magnitudes have not to one another the ratio which
a number has to a number.
\end{claim}
\begin{evidence}[Proof of X.7]
\label{ev:X.7}
Contrapositive of X.6.
\dependson{X.7}{X.6}
\end{evidence}

\begin{claim}[Proposition X.8: No number ratio implies incommensurable]
\label{prop:X.8}
If two magnitudes have not to one another the ratio which a number
has to a number, the magnitudes will be incommensurable.
\end{claim}
\begin{evidence}[Proof of X.8]
\label{ev:X.8}
Contrapositive of X.5.
\dependson{X.8}{X.5}
\end{evidence}

\begin{claim}[Proposition X.9: Commensurability of squares from commensurability of sides]
\label{prop:X.9}
The squares on straight lines commensurable in length have to one
another the ratio which a square number has to a square number; and
squares which have to one another the ratio which a square number
has to a square number will also have their sides commensurable in
length.
\end{claim}
\begin{evidence}[Proof of X.9]
\label{ev:X.9}
If $a : b = m : n$ (integers) then $a^2 : b^2 = m^2 : n^2$; converse
holds by VIII.14 applied to integers and VI.22 applied to magnitudes.
\dependson{X.9}{VI.22}
\dependson{X.9}{VIII.14}
\dependson{X.9}{X.5}
\dependson{X.9}{X.6}
\end{evidence}

\begin{claim}[Proposition X.10: Construct incommensurables]
\label{prop:X.10}
To find two straight lines incommensurable, the one in length only,
the other in square also, with an assigned straight line.
\end{claim}
\begin{evidence}[Proof of X.10]
\label{ev:X.10}
Take a non-square integer ratio (e.g.\ $2 : 1$) and use it via X.9 to
construct an incommensurable-in-length pair; build the
incommensurable-in-square one similarly with a ratio that is neither
square nor cube.
\dependson{X.10}{VI.13}
\dependson{X.10}{X.9}
\end{evidence}

\begin{claim}[Proposition X.11: Commensurability is transitive]
\label{prop:X.11}
If four magnitudes be proportional, and the first be commensurable
with the second, the third also will be commensurable with the
fourth; and if the first be incommensurable with the second, the
third also will be incommensurable with the fourth.
\end{claim}
\begin{evidence}[Proof of X.11]
\label{ev:X.11}
Commensurability $\iff$ rational ratio (X.5 / X.6); rational ratios
are preserved under equality of ratios.
\dependson{X.11}{X.5}
\dependson{X.11}{X.6}
\end{evidence}

\begin{claim}[Proposition X.12: Commensurability is transitive (three magnitudes)]
\label{prop:X.12}
Magnitudes commensurable with the same magnitude are commensurable
with one another.
\end{claim}
\begin{evidence}[Proof of X.12]
\label{ev:X.12}
If $a \sim c$ and $b \sim c$ (commensurable), then $a \sim b$ by
composing ratios (X.11).
\dependson{X.12}{X.11}
\end{evidence}

\begin{claim}[Proposition X.13: Incommensurable preserved through transitivity]
\label{prop:X.13}
If two magnitudes be commensurable, and one of them be incommensurable
with any magnitude, the remaining one will also be incommensurable
with the same.
\end{claim}
\begin{evidence}[Proof of X.13]
\label{ev:X.13}
Contrapositive of X.12.
\dependson{X.13}{X.12}
\end{evidence}

\begin{claim}[Proposition X.14: Squares preserve commensurability of sides]
\label{prop:X.14}
If four straight lines be proportional, and the square on the first
be greater than the square on the second by the square on a straight
line commensurable with the first, the square on the third will also
be greater than the square on the fourth by the square on a straight
line commensurable with the third.
\end{claim}
\begin{evidence}[Proof of X.14]
\label{ev:X.14}
Proportionality lifts to the squares; the deviation magnitude
inherits the commensurability relation.
\dependson{X.14}{VI.22}
\dependson{X.14}{X.11}
\end{evidence}

\begin{claim}[Proposition X.15: Sum of commensurables is commensurable]
\label{prop:X.15}
If two commensurable magnitudes be added together, the whole will
also be commensurable with each of them; and if the whole be
commensurable with one of them, the original magnitudes will also be
commensurable.
\end{claim}
\begin{evidence}[Proof of X.15]
\label{ev:X.15}
$a, b$ have a common measure $d$, so $a + b = (m + n) d$ shares $d$.
Converse: if $a + b$ and $a$ share a measure, then so does $b$ by
Common Notion 3.
\dependson{X.15}{cn:3}
\dependson{X.15}{def:X.1}
\end{evidence}

\begin{claim}[Proposition X.16: Sum with incommensurable]
\label{prop:X.16}
If two incommensurable magnitudes be added together, the whole will
also be incommensurable with each of them; and if the whole be
incommensurable with one of them, the original magnitudes will also
be incommensurable.
\end{claim}
\begin{evidence}[Proof of X.16]
\label{ev:X.16}
Contrapositive of X.15.
\dependson{X.16}{X.15}
\end{evidence}

\begin{claim}[Proposition X.17: Application of areas with commensurable difference]
\label{prop:X.17}
If there be two unequal straight lines, and to the greater there be
applied a parallelogram equal to the fourth part of the square on the
less and deficient by a square figure, and if it divide it into parts
which are commensurable in length, then the square on the greater
will be greater than the square on the less by the square on a
straight line commensurable in length with the greater.
\end{claim}
\begin{evidence}[Proof of X.17]
\label{ev:X.17}
By VI.28 / VI.29 the application of areas with deficient/excess
square corresponds to solving a quadratic; commensurability of the
parts forces commensurability of the discriminant.
\dependson{X.17}{VI.28}
\dependson{X.17}{X.14}
\dependson{X.17}{X.15}
\end{evidence}

\begin{claim}[Proposition X.18: Application of areas with incommensurable difference]
\label{prop:X.18}
If there be two unequal straight lines, and to the greater there be
applied a parallelogram equal to the fourth part of the square on the
less and deficient by a square figure, and if it divide it into parts
incommensurable in length, then the square on the greater will be
greater than the square on the less by the square on a straight line
incommensurable in length with the greater.
\end{claim}
\begin{evidence}[Proof of X.18]
\label{ev:X.18}
Same construction as X.17 with the opposite hypothesis;
incommensurability of the parts forces incommensurability of the
discriminant.
\dependson{X.18}{VI.28}
\dependson{X.18}{X.16}
\dependson{X.18}{X.17}
\end{evidence}

\begin{claim}[Proposition X.19: Rectangle on rationals is rational]
\label{prop:X.19}
The rectangle contained by rational straight lines commensurable in
length is rational.
\end{claim}
\begin{evidence}[Proof of X.19]
\label{ev:X.19}
Rational sides commensurable in length have integer ratios; product
of two such sides is in rational ratio to the assigned-square area.
\dependson{X.19}{X.9}
\dependson{X.19}{def:X.3}
\dependson{X.19}{def:X.4}
\end{evidence}

\begin{claim}[Proposition X.20: Rational area divided by rational side is rational]
\label{prop:X.20}
If a rational area be applied to a rational straight line, it produces
as breadth a straight line rational and commensurable in length with
the straight line to which it is applied.
\end{claim}
\begin{evidence}[Proof of X.20]
\label{ev:X.20}
$A = \ell \cdot b$ with $A$ and $\ell$ rational forces $b$ rational by
X.19 / X.9.
\dependson{X.20}{X.19}
\end{evidence}

\begin{claim}[Proposition X.21: Medial rectangle]
\label{prop:X.21}
The rectangle contained by rational straight lines commensurable in
square only is irrational, and the side of the square equal to it is
irrational.  Let the latter be called medial.
\end{claim}
\begin{evidence}[Proof of X.21]
\label{ev:X.21}
Commensurable-in-square-only means the square on each is rational
but the lengths are not in integer ratio.  The rectangle is then in
a non-rational ratio to a rational area; its square root is the
medial straight line (Definition XIII.3).
\dependson{X.21}{X.9}
\dependson{X.21}{def:X.3}
\dependson{X.21}{def:XIII.3}
\end{evidence}

\begin{claim}[Proposition X.22: Square on a medial]
\label{prop:X.22}
The square on a medial straight line, if applied to a rational
straight line, produces as breadth a straight line rational and
incommensurable in length with that to which it is applied.
\end{claim}
\begin{evidence}[Proof of X.22]
\label{ev:X.22}
A medial squared is rational; applying it to a rational base, the
breadth is rational; the incommensurability follows from the fact
that the medial is incommensurable in length with the rational.
\dependson{X.22}{X.21}
\end{evidence}

\begin{claim}[Proposition X.23: Magnitudes commensurable with medial are medial]
\label{prop:X.23}
A straight line commensurable with a medial straight line is medial.
\end{claim}
\begin{evidence}[Proof of X.23]
\label{ev:X.23}
Commensurability preserves the medial property: scaling a medial by
a rational ratio leaves it medial.
\dependson{X.23}{X.21}
\dependson{X.23}{def:XIII.3}
\end{evidence}

\begin{claim}[Proposition X.24: Rectangle of commensurable medials]
\label{prop:X.24}
The rectangle contained by medial straight lines commensurable in
length is medial.
\end{claim}
\begin{evidence}[Proof of X.24]
\label{ev:X.24}
Product of two medials in rational length-ratio is again the
geometric mean of two rationals (Definition XIII.3).
\dependson{X.24}{X.21}
\dependson{X.24}{X.23}
\end{evidence}

\begin{claim}[Proposition X.25: Rectangle of medials commensurable in square only]
\label{prop:X.25}
The rectangle contained by medial straight lines commensurable in
square only is either rational or medial.
\end{claim}
\begin{evidence}[Proof of X.25]
\label{ev:X.25}
Two cases depending on whether the rectangle has a rational
square-root.  Both cases are realised by explicit constructions.
\dependson{X.25}{X.21}
\dependson{X.25}{X.24}
\end{evidence}

\begin{claim}[Proposition X.26: Medial difference is not rational]
\label{prop:X.26}
A medial area does not exceed a medial area by a rational area.
\end{claim}
\begin{evidence}[Proof of X.26]
\label{ev:X.26}
A difference $M_1 - M_2 = R$ with $R$ rational and $M_1$, $M_2$
medial would force $M_1$, $M_2$ to be in a rational ratio,
contradicting their being medial in distinct families.
\dependson{X.26}{X.21}
\dependson{X.26}{X.25}
\end{evidence}

\begin{claim}[Proposition X.27: Two medial lines commensurable in square]
\label{prop:X.27}
To find medial straight lines commensurable in square only which
contain a rational rectangle.
\end{claim}
\begin{evidence}[Proof of X.27]
\label{ev:X.27}
Construct two such medials from a fixed rational by extracting two
square-roots of ratios in lowest terms.
\dependson{X.27}{X.21}
\dependson{X.27}{X.22}
\end{evidence}

\begin{claim}[Proposition X.28: Medials enclosing a medial rectangle]
\label{prop:X.28}
To find medial straight lines commensurable in square only which
contain a medial rectangle.
\end{claim}
\begin{evidence}[Proof of X.28]
\label{ev:X.28}
Same construction as X.27 with an extra medial step.
\dependson{X.28}{X.21}
\dependson{X.28}{X.27}
\end{evidence}

\begin{claim}[Proposition X.29: Two rationals commensurable in square, square-difference of commensurable kind, Lemma 1]
\label{prop:X.29}
To find two rational straight lines commensurable in square only such
that the square on the greater is greater than the square on the
less by the square on a straight line commensurable in length with
the greater.
\end{claim}
\begin{evidence}[Proof of X.29]
\label{ev:X.29}
Take a rational line $a$ and apply X.17 with a deficient square; the
construction gives the required pair.
\dependson{X.29}{X.17}
\end{evidence}

\begin{claim}[Proposition X.30: Same as X.29 with the discriminant incommensurable]
\label{prop:X.30}
To find two rational straight lines commensurable in square only such
that the square on the greater is greater than the square on the
less by the square on a straight line incommensurable in length with
the greater.
\end{claim}
\begin{evidence}[Proof of X.30]
\label{ev:X.30}
Apply X.18 (the incommensurable analogue of X.17).
\dependson{X.30}{X.18}
\end{evidence}

\begin{claim}[Proposition X.31: Two medials with rational discriminant relation]
\label{prop:X.31}
To find two medial straight lines commensurable in square only,
containing a rational rectangle, such that the square on the greater
is greater than the square on the less by the square on a straight
line commensurable in length with the greater.
\end{claim}
\begin{evidence}[Proof of X.31]
\label{ev:X.31}
Combine X.27 with the discriminant-control of X.29.
\dependson{X.31}{X.27}
\dependson{X.31}{X.29}
\end{evidence}

\begin{claim}[Proposition X.32: Two medials, medial rectangle, commensurable discriminant]
\label{prop:X.32}
To find two medial straight lines commensurable in square only,
containing a medial rectangle, such that the square on the greater
is greater than the square on the less by the square on a straight
line commensurable in length with the greater.
\end{claim}
\begin{evidence}[Proof of X.32]
\label{ev:X.32}
Same as X.31 with X.28 in place of X.27.
\dependson{X.32}{X.28}
\dependson{X.32}{X.31}
\end{evidence}

\begin{claim}[Proposition X.33: Sum of squares rational, rectangle medial, sides incommensurable in square]
\label{prop:X.33}
To find two straight lines incommensurable in square which make the
sum of the squares on them rational but the rectangle contained by
them medial.
\end{claim}
\begin{evidence}[Proof of X.33]
\label{ev:X.33}
Apply X.30: a difference-of-squares construction with an
incommensurable discriminant produces such a pair.
\dependson{X.33}{X.30}
\end{evidence}

\begin{claim}[Proposition X.34: Sum medial, rectangle rational]
\label{prop:X.34}
To find two straight lines incommensurable in square which make the
sum of the squares on them medial but the rectangle contained by
them rational.
\end{claim}
\begin{evidence}[Proof of X.34]
\label{ev:X.34}
Variant of X.33 with the medial/rational roles swapped, via X.31.
\dependson{X.34}{X.31}
\dependson{X.34}{X.33}
\end{evidence}

\begin{claim}[Proposition X.35: Sum medial, rectangle medial, sum incommensurable with rectangle]
\label{prop:X.35}
To find two straight lines incommensurable in square which make the
sum of the squares on them medial and the rectangle contained by
them medial and moreover incommensurable with the sum of the squares
on them.
\end{claim}
\begin{evidence}[Proof of X.35]
\label{ev:X.35}
Variant of X.34 with both quantities medial; use X.32.
\dependson{X.35}{X.32}
\dependson{X.35}{X.34}
\end{evidence}

\begin{claim}[Proposition X.36: Binomial straight line]
\label{prop:X.36}
If two rational straight lines commensurable in square only be added
together, the whole is irrational; and let it be called binomial.
\end{claim}
\begin{evidence}[Proof of X.36]
\label{ev:X.36}
The sum $a + b$ with $a$, $b$ rational and incommensurable in length
has square $a^2 + b^2 + 2ab$ where $2ab$ is medial (X.21), so the
square on the sum is the sum of a rational and a medial: irrational.
\dependson{X.36}{X.21}
\dependson{X.36}{def:X.II.1}
\end{evidence}

\begin{claim}[Proposition X.37: First bimedial straight line]
\label{prop:X.37}
If two medial straight lines commensurable in square only and
containing a rational rectangle be added together, the whole is
irrational; and let it be called first bimedial.
\end{claim}
\begin{evidence}[Proof of X.37]
\label{ev:X.37}
Sum of two medials with rational rectangle; the square consists of
two medials and a rational --- irrational.
\dependson{X.37}{X.27}
\dependson{X.37}{X.36}
\end{evidence}

\begin{claim}[Proposition X.38: Second bimedial straight line]
\label{prop:X.38}
If two medial straight lines commensurable in square only and
containing a medial rectangle be added together, the whole is
irrational; and let it be called second bimedial.
\end{claim}
\begin{evidence}[Proof of X.38]
\label{ev:X.38}
Same scheme as X.37 with the rectangle medial instead of rational.
\dependson{X.38}{X.28}
\dependson{X.38}{X.37}
\end{evidence}

\begin{claim}[Proposition X.39: Major straight line]
\label{prop:X.39}
If two straight lines incommensurable in square which make the sum
of the squares on them rational, but the rectangle contained by them
medial, be added together, the whole straight line is irrational;
and let it be called major.
\end{claim}
\begin{evidence}[Proof of X.39]
\label{ev:X.39}
Sum of an X.33 pair has square = rational + medial: irrational.
\dependson{X.39}{X.33}
\dependson{X.39}{X.36}
\end{evidence}

\begin{claim}[Proposition X.40: Side of a rational plus medial area]
\label{prop:X.40}
If two straight lines incommensurable in square which make the sum
of the squares on them medial, but the rectangle contained by them
rational, be added together, the whole straight line is irrational;
and let it be called the side of a rational plus a medial area.
\end{claim}
\begin{evidence}[Proof of X.40]
\label{ev:X.40}
Sum of an X.34 pair; same scheme as X.39.
\dependson{X.40}{X.34}
\dependson{X.40}{X.39}
\end{evidence}

\begin{claim}[Proposition X.41: Side of the sum of two medial areas]
\label{prop:X.41}
If two straight lines incommensurable in square which make the sum
of the squares on them medial, and the rectangle contained by them
medial and also incommensurable with the sum of the squares on them,
be added together, the remaining straight line is irrational; and
let it be called the side of the sum of two medial areas.
\end{claim}
\begin{evidence}[Proof of X.41]
\label{ev:X.41}
Sum of an X.35 pair; same scheme as X.39.
\dependson{X.41}{X.35}
\dependson{X.41}{X.40}
\end{evidence}

\begin{claim}[Proposition X.42: A binomial has unique decomposition]
\label{prop:X.42}
A binomial straight line is divided into its terms at one point only.
\end{claim}
\begin{evidence}[Proof of X.42]
\label{ev:X.42}
Suppose two decompositions $a_1 + b_1 = a_2 + b_2$ of the same
binomial.  Comparing rationals and medials in the squares forces
$a_1 = a_2$ and $b_1 = b_2$.
\dependson{X.42}{X.36}
\end{evidence}

\begin{claim}[Proposition X.43: A first bimedial has unique decomposition]
\label{prop:X.43}
A first bimedial straight line is divided at one and the same point
only.
\end{claim}
\begin{evidence}[Proof of X.43]
\label{ev:X.43}
Same uniqueness argument as X.42 applied to the first bimedial.
\dependson{X.43}{X.37}
\dependson{X.43}{X.42}
\end{evidence}

\begin{claim}[Proposition X.44: A second bimedial has unique decomposition]
\label{prop:X.44}
A second bimedial straight line is divided at one point only.
\end{claim}
\begin{evidence}[Proof of X.44]
\label{ev:X.44}
Same uniqueness argument applied to the second bimedial.
\dependson{X.44}{X.38}
\dependson{X.44}{X.43}
\end{evidence}

\begin{claim}[Proposition X.45: A major has unique decomposition]
\label{prop:X.45}
A major straight line is divided at one and the same point only.
\end{claim}
\begin{evidence}[Proof of X.45]
\label{ev:X.45}
Same uniqueness argument applied to the major.
\dependson{X.45}{X.39}
\dependson{X.45}{X.44}
\end{evidence}

\begin{claim}[Proposition X.46: Side of rational+medial has unique decomposition]
\label{prop:X.46}
The side of a rational plus a medial area is divided at one and the
same point only.
\end{claim}
\begin{evidence}[Proof of X.46]
\label{ev:X.46}
Same uniqueness argument.
\dependson{X.46}{X.40}
\dependson{X.46}{X.45}
\end{evidence}

\begin{claim}[Proposition X.47: Side of two medial areas has unique decomposition]
\label{prop:X.47}
The side of the sum of two medial areas is divided at one and the
same point only.
\end{claim}
\begin{evidence}[Proof of X.47]
\label{ev:X.47}
Same uniqueness argument.
\dependson{X.47}{X.41}
\dependson{X.47}{X.46}
\end{evidence}

\begin{claim}[Proposition X.48: First binomial straight line]
\label{prop:X.48}
To find the first binomial straight line.
\end{claim}
\begin{evidence}[Proof of X.48]
\label{ev:X.48}
Construct $a + b$ with $a$ commensurable in length with the
assigned rational and the square-discriminant commensurable with the
greater (X.29).
\dependson{X.48}{X.29}
\dependson{X.48}{X.36}
\dependson{X.48}{def:X.II.1}
\end{evidence}

\begin{claim}[Proposition X.49: Second binomial straight line]
\label{prop:X.49}
To find the second binomial straight line.
\end{claim}
\begin{evidence}[Proof of X.49]
\label{ev:X.49}
Construct as in X.48 but with $b$ (rather than $a$) commensurable
with the assigned rational.
\dependson{X.49}{X.48}
\dependson{X.49}{def:X.II.2}
\end{evidence}

\begin{claim}[Proposition X.50: Third binomial straight line]
\label{prop:X.50}
To find the third binomial straight line.
\end{claim}
\begin{evidence}[Proof of X.50]
\label{ev:X.50}
Neither term commensurable with the assigned rational, but the
square-discriminant commensurable with the greater (X.29 variant).
\dependson{X.50}{X.29}
\dependson{X.50}{X.48}
\dependson{X.50}{def:X.II.3}
\end{evidence}

\begin{claim}[Proposition X.51: Fourth binomial straight line]
\label{prop:X.51}
To find the fourth binomial straight line.
\end{claim}
\begin{evidence}[Proof of X.51]
\label{ev:X.51}
$a$ commensurable with the assigned rational, square-discriminant
incommensurable with $a$ (X.30).
\dependson{X.51}{X.30}
\dependson{X.51}{X.48}
\dependson{X.51}{def:X.II.4}
\end{evidence}

\begin{claim}[Proposition X.52: Fifth binomial straight line]
\label{prop:X.52}
To find the fifth binomial straight line.
\end{claim}
\begin{evidence}[Proof of X.52]
\label{ev:X.52}
$b$ commensurable, discriminant incommensurable with the greater.
\dependson{X.52}{X.30}
\dependson{X.52}{X.51}
\dependson{X.52}{def:X.II.5}
\end{evidence}

\begin{claim}[Proposition X.53: Sixth binomial straight line]
\label{prop:X.53}
To find the sixth binomial straight line.
\end{claim}
\begin{evidence}[Proof of X.53]
\label{ev:X.53}
Neither term commensurable, discriminant incommensurable.
\dependson{X.53}{X.30}
\dependson{X.53}{X.52}
\dependson{X.53}{def:X.II.6}
\end{evidence}

\begin{claim}[Proposition X.54: Rectangle on a first binomial is rational on rational]
\label{prop:X.54}
If an area be contained by a rational straight line and the first
binomial, the side of the area is the irrational straight line which
is called binomial.
\end{claim}
\begin{evidence}[Proof of X.54]
\label{ev:X.54}
$\sqrt{R \cdot \text{first binomial}}$ has the form of a binomial in
the assigned rational base.
\dependson{X.54}{X.36}
\dependson{X.54}{X.48}
\end{evidence}

\begin{claim}[Proposition X.55: Rectangle on a second binomial yields a first bimedial]
\label{prop:X.55}
If an area be contained by a rational straight line and the second
binomial, the side of the area is the irrational straight line which
is called first bimedial.
\end{claim}
\begin{evidence}[Proof of X.55]
\label{ev:X.55}
Same pattern as X.54.
\dependson{X.55}{X.37}
\dependson{X.55}{X.49}
\dependson{X.55}{X.54}
\end{evidence}

\begin{claim}[Proposition X.56: Rectangle on a third binomial yields a second bimedial]
\label{prop:X.56}
If an area be contained by a rational straight line and the third
binomial, the side of the area is the irrational straight line which
is called second bimedial.
\end{claim}
\begin{evidence}[Proof of X.56]
\label{ev:X.56}
Same pattern.
\dependson{X.56}{X.38}
\dependson{X.56}{X.50}
\dependson{X.56}{X.55}
\end{evidence}

\begin{claim}[Proposition X.57: Rectangle on a fourth binomial yields a major]
\label{prop:X.57}
If an area be contained by a rational straight line and the fourth
binomial, the side of the area is the irrational straight line which
is called major.
\end{claim}
\begin{evidence}[Proof of X.57]
\label{ev:X.57}
Same pattern.
\dependson{X.57}{X.39}
\dependson{X.57}{X.51}
\dependson{X.57}{X.56}
\end{evidence}

\begin{claim}[Proposition X.58: Rectangle on a fifth binomial yields a side of rational+medial]
\label{prop:X.58}
If an area be contained by a rational straight line and the fifth
binomial, the side of the area is the irrational straight line which
is the side of a rational plus a medial area.
\end{claim}
\begin{evidence}[Proof of X.58]
\label{ev:X.58}
Same pattern.
\dependson{X.58}{X.40}
\dependson{X.58}{X.52}
\dependson{X.58}{X.57}
\end{evidence}

\begin{claim}[Proposition X.59: Rectangle on a sixth binomial yields side of two medials]
\label{prop:X.59}
If an area be contained by a rational straight line and the sixth
binomial, the side of the area is the irrational straight line which
is called the side of the sum of two medial areas.
\end{claim}
\begin{evidence}[Proof of X.59]
\label{ev:X.59}
Same pattern.
\dependson{X.59}{X.41}
\dependson{X.59}{X.53}
\dependson{X.59}{X.58}
\end{evidence}

\begin{claim}[Proposition X.60: Square on a binomial yields a first binomial]
\label{prop:X.60}
The square on the binomial straight line applied to a rational
straight line produces as breadth the first binomial.
\end{claim}
\begin{evidence}[Proof of X.60]
\label{ev:X.60}
Inverse of X.54: squaring and dividing by the assigned rational
recovers the first binomial.
\dependson{X.60}{X.48}
\dependson{X.60}{X.54}
\end{evidence}

\begin{claim}[Proposition X.61: Square on a first bimedial yields a second binomial]
\label{prop:X.61}
The square on the first bimedial straight line applied to a rational
straight line produces as breadth the second binomial.
\end{claim}
\begin{evidence}[Proof of X.61]
\label{ev:X.61}
Inverse of X.55.
\dependson{X.61}{X.49}
\dependson{X.61}{X.55}
\dependson{X.61}{X.60}
\end{evidence}

\begin{claim}[Proposition X.62: Square on a second bimedial yields a third binomial]
\label{prop:X.62}
The square on the second bimedial straight line applied to a rational
straight line produces as breadth the third binomial.
\end{claim}
\begin{evidence}[Proof of X.62]
\label{ev:X.62}
Inverse of X.56.
\dependson{X.62}{X.50}
\dependson{X.62}{X.56}
\dependson{X.62}{X.61}
\end{evidence}

\begin{claim}[Proposition X.63: Square on a major yields a fourth binomial]
\label{prop:X.63}
The square on the major straight line applied to a rational straight
line produces as breadth the fourth binomial.
\end{claim}
\begin{evidence}[Proof of X.63]
\label{ev:X.63}
Inverse of X.57.
\dependson{X.63}{X.51}
\dependson{X.63}{X.57}
\dependson{X.63}{X.62}
\end{evidence}

\begin{claim}[Proposition X.64: Square on side-of-rational-plus-medial yields a fifth binomial]
\label{prop:X.64}
The square on the side of a rational plus a medial area applied to
a rational straight line produces as breadth the fifth binomial.
\end{claim}
\begin{evidence}[Proof of X.64]
\label{ev:X.64}
Inverse of X.58.
\dependson{X.64}{X.52}
\dependson{X.64}{X.58}
\dependson{X.64}{X.63}
\end{evidence}

\begin{claim}[Proposition X.65: Square on side-of-two-medials yields a sixth binomial]
\label{prop:X.65}
The square on the side of the sum of two medial areas applied to a
rational straight line produces as breadth the sixth binomial.
\end{claim}
\begin{evidence}[Proof of X.65]
\label{ev:X.65}
Inverse of X.59.
\dependson{X.65}{X.53}
\dependson{X.65}{X.59}
\dependson{X.65}{X.64}
\end{evidence}

\begin{claim}[Proposition X.66: Commensurable with binomial is binomial]
\label{prop:X.66}
A straight line commensurable in length with a binomial straight
line is itself also binomial and the same in order.
\end{claim}
\begin{evidence}[Proof of X.66]
\label{ev:X.66}
Multiplication by a rational ratio preserves the binomial type and
order.
\dependson{X.66}{X.36}
\dependson{X.66}{X.48}
\end{evidence}

\begin{claim}[Proposition X.67: Commensurable with bimedial is bimedial]
\label{prop:X.67}
A straight line commensurable in length with a bimedial straight
line is itself bimedial and the same in order.
\end{claim}
\begin{evidence}[Proof of X.67]
\label{ev:X.67}
Same scheme as X.66.
\dependson{X.67}{X.37}
\dependson{X.67}{X.38}
\dependson{X.67}{X.66}
\end{evidence}

\begin{claim}[Proposition X.68: Commensurable with major is major]
\label{prop:X.68}
A straight line commensurable with a major straight line is itself
major.
\end{claim}
\begin{evidence}[Proof of X.68]
\label{ev:X.68}
Same scheme.
\dependson{X.68}{X.39}
\dependson{X.68}{X.67}
\end{evidence}

\begin{claim}[Proposition X.69: Commensurable with side of rational+medial is the same]
\label{prop:X.69}
A straight line commensurable with the side of a rational plus a
medial area is itself such a side.
\end{claim}
\begin{evidence}[Proof of X.69]
\label{ev:X.69}
Same scheme.
\dependson{X.69}{X.40}
\dependson{X.69}{X.68}
\end{evidence}

\begin{claim}[Proposition X.70: Commensurable with side of two-medials is the same]
\label{prop:X.70}
A straight line commensurable with the side of the sum of two medial
areas is itself such a side.
\end{claim}
\begin{evidence}[Proof of X.70]
\label{ev:X.70}
Same scheme.
\dependson{X.70}{X.41}
\dependson{X.70}{X.69}
\end{evidence}

\begin{claim}[Proposition X.71: Rational + medial sum is one of the four irrationals]
\label{prop:X.71}
If a rational and a medial area be added together, four irrational
straight lines arise, namely either a binomial, a first bimedial, a
major, or a side of a rational plus a medial area.
\end{claim}
\begin{evidence}[Proof of X.71]
\label{ev:X.71}
The square of any of the four classes (X.36, X.37, X.39, X.40) is
the sum of a rational and a medial; conversely, every such sum
arises in exactly one of these forms.
\dependson{X.71}{X.36}
\dependson{X.71}{X.37}
\dependson{X.71}{X.39}
\dependson{X.71}{X.40}
\end{evidence}

\begin{claim}[Proposition X.72: Medial + medial sum yields a bimedial or a side of two medials]
\label{prop:X.72}
If two medial areas incommensurable with one another be added
together, the remaining two irrational straight lines arise, namely
either a second bimedial or a side of the sum of two medial areas.
\end{claim}
\begin{evidence}[Proof of X.72]
\label{ev:X.72}
The square of X.38 or X.41 is a sum of two incommensurable medial
areas; converse runs the same way.
\dependson{X.72}{X.38}
\dependson{X.72}{X.41}
\dependson{X.72}{X.71}
\end{evidence}

\begin{claim}[Proposition X.73: Apotome straight line]
\label{prop:X.73}
If from a rational straight line there be subtracted a rational
straight line commensurable with the whole in square only, the
remainder is irrational; and let it be called apotome.
\end{claim}
\begin{evidence}[Proof of X.73]
\label{ev:X.73}
$a - b$ with $a$, $b$ commensurable in square only is the negation
of the binomial case (X.36); the same argument shows it is
irrational.
\dependson{X.73}{X.36}
\end{evidence}

\begin{claim}[Proposition X.74: First apotome of a medial]
\label{prop:X.74}
If from a medial straight line there be subtracted a medial straight
line commensurable with the whole in square only, and containing
with the whole a rational rectangle, the remainder is irrational;
and let it be called first apotome of a medial.
\end{claim}
\begin{evidence}[Proof of X.74]
\label{ev:X.74}
Negation of X.37.
\dependson{X.74}{X.37}
\dependson{X.74}{X.73}
\end{evidence}

\begin{claim}[Proposition X.75: Second apotome of a medial]
\label{prop:X.75}
If from a medial straight line there be subtracted a medial straight
line commensurable with the whole in square only, and containing
with the whole a medial rectangle, the remainder is irrational; and
let it be called second apotome of a medial.
\end{claim}
\begin{evidence}[Proof of X.75]
\label{ev:X.75}
Negation of X.38.
\dependson{X.75}{X.38}
\dependson{X.75}{X.74}
\end{evidence}

\begin{claim}[Proposition X.76: Minor straight line]
\label{prop:X.76}
If from a straight line there be subtracted a straight line
incommensurable in square with the whole, which with the whole makes
the squares on them added together rational, but the rectangle
contained by them medial, the remainder is irrational; and let it
be called minor.
\end{claim}
\begin{evidence}[Proof of X.76]
\label{ev:X.76}
Negation of X.39.  This is the "minor" line (Definition XIII.4).
\dependson{X.76}{X.39}
\dependson{X.76}{def:XIII.4}
\end{evidence}

\begin{claim}[Proposition X.77: Line producing with rational area a medial whole]
\label{prop:X.77}
If from a straight line there be subtracted a straight line
incommensurable in square with the whole which with the whole makes
the sum of squares medial but twice the rectangle rational, the
remainder is irrational; let it be called that which produces with
a rational area a medial whole.
\end{claim}
\begin{evidence}[Proof of X.77]
\label{ev:X.77}
Negation of X.40.
\dependson{X.77}{X.40}
\dependson{X.77}{X.76}
\dependson{X.77}{def:XIII.5}
\end{evidence}

\begin{claim}[Proposition X.78: Line producing with medial area a medial whole]
\label{prop:X.78}
If from a straight line there be subtracted a straight line
incommensurable in square with the whole which with the whole makes
both the sum of squares and twice the rectangle medial and the two
sums incommensurable with one another, the remainder is irrational;
let it be called that which produces with a medial area a medial
whole.
\end{claim}
\begin{evidence}[Proof of X.78]
\label{ev:X.78}
Negation of X.41.
\dependson{X.78}{X.41}
\dependson{X.78}{X.77}
\end{evidence}

\begin{claim}[Proposition X.79: Apotome has unique annex]
\label{prop:X.79}
Only one rational straight line can be annexed to an apotome which
is commensurable with the whole in square only.
\end{claim}
\begin{evidence}[Proof of X.79]
\label{ev:X.79}
Uniqueness analogue of X.42 for apotomes.
\dependson{X.79}{X.42}
\dependson{X.79}{X.73}
\end{evidence}

\begin{claim}[Proposition X.80: First-apotome-of-medial uniqueness]
\label{prop:X.80}
Only one medial straight line can be annexed to a first apotome of
a medial which is commensurable with the whole in square only and
forms with it a rational rectangle.
\end{claim}
\begin{evidence}[Proof of X.80]
\label{ev:X.80}
Same uniqueness pattern.
\dependson{X.80}{X.43}
\dependson{X.80}{X.74}
\end{evidence}

\begin{claim}[Proposition X.81: Second-apotome-of-medial uniqueness]
\label{prop:X.81}
Only one medial straight line can be annexed to a second apotome of
a medial which is commensurable with the whole in square only and
forms with it a medial rectangle.
\end{claim}
\begin{evidence}[Proof of X.81]
\label{ev:X.81}
Same uniqueness pattern.
\dependson{X.81}{X.44}
\dependson{X.81}{X.75}
\end{evidence}

\begin{claim}[Proposition X.82: Minor uniqueness]
\label{prop:X.82}
Only one straight line can be annexed to a minor.
\end{claim}
\begin{evidence}[Proof of X.82]
\label{ev:X.82}
Same uniqueness pattern.
\dependson{X.82}{X.45}
\dependson{X.82}{X.76}
\end{evidence}

\begin{claim}[Proposition X.83: Uniqueness for X.77's line]
\label{prop:X.83}
Only one straight line can be annexed to the line producing with a
rational area a medial whole.
\end{claim}
\begin{evidence}[Proof of X.83]
\label{ev:X.83}
Same uniqueness pattern.
\dependson{X.83}{X.46}
\dependson{X.83}{X.77}
\end{evidence}

\begin{claim}[Proposition X.84: Uniqueness for X.78's line]
\label{prop:X.84}
Only one straight line can be annexed to the line producing with a
medial area a medial whole.
\end{claim}
\begin{evidence}[Proof of X.84]
\label{ev:X.84}
Same uniqueness pattern.
\dependson{X.84}{X.47}
\dependson{X.84}{X.78}
\end{evidence}

\begin{claim}[Proposition X.85: First apotome]
\label{prop:X.85}
To find the first apotome.
\end{claim}
\begin{evidence}[Proof of X.85]
\label{ev:X.85}
Take a first binomial $a + b$ (X.48); the difference $a - b$ is the
first apotome.
\dependson{X.85}{X.48}
\dependson{X.85}{X.73}
\dependson{X.85}{def:X.III.1}
\end{evidence}

\begin{claim}[Proposition X.86: Second apotome]
\label{prop:X.86}
To find the second apotome.
\end{claim}
\begin{evidence}[Proof of X.86]
\label{ev:X.86}
Use the second binomial as the model (X.49).
\dependson{X.86}{X.49}
\dependson{X.86}{X.85}
\dependson{X.86}{def:X.III.2}
\end{evidence}

\begin{claim}[Proposition X.87: Third apotome]
\label{prop:X.87}
To find the third apotome.
\end{claim}
\begin{evidence}[Proof of X.87]
\label{ev:X.87}
Use the third binomial (X.50).
\dependson{X.87}{X.50}
\dependson{X.87}{X.86}
\dependson{X.87}{def:X.III.3}
\end{evidence}

\begin{claim}[Proposition X.88: Fourth apotome]
\label{prop:X.88}
To find the fourth apotome.
\end{claim}
\begin{evidence}[Proof of X.88]
\label{ev:X.88}
Use the fourth binomial (X.51).
\dependson{X.88}{X.51}
\dependson{X.88}{X.87}
\dependson{X.88}{def:X.III.4}
\end{evidence}

\begin{claim}[Proposition X.89: Fifth apotome]
\label{prop:X.89}
To find the fifth apotome.
\end{claim}
\begin{evidence}[Proof of X.89]
\label{ev:X.89}
Use the fifth binomial (X.52).
\dependson{X.89}{X.52}
\dependson{X.89}{X.88}
\dependson{X.89}{def:X.III.5}
\end{evidence}

\begin{claim}[Proposition X.90: Sixth apotome]
\label{prop:X.90}
To find the sixth apotome.
\end{claim}
\begin{evidence}[Proof of X.90]
\label{ev:X.90}
Use the sixth binomial (X.53).
\dependson{X.90}{X.53}
\dependson{X.90}{X.89}
\dependson{X.90}{def:X.III.6}
\end{evidence}

\begin{claim}[Proposition X.91: Side of first-apotome area is an apotome]
\label{prop:X.91}
If an area be contained by a rational straight line and a first
apotome, the side of the area is an apotome.
\end{claim}
\begin{evidence}[Proof of X.91]
\label{ev:X.91}
Dual of X.54 for apotomes.
\dependson{X.91}{X.54}
\dependson{X.91}{X.85}
\end{evidence}

\begin{claim}[Proposition X.92: Side of second-apotome area is a first apotome of medial]
\label{prop:X.92}
If an area be contained by a rational straight line and a second
apotome, the side of the area is a first apotome of a medial.
\end{claim}
\begin{evidence}[Proof of X.92]
\label{ev:X.92}
Dual of X.55.
\dependson{X.92}{X.55}
\dependson{X.92}{X.86}
\dependson{X.92}{X.91}
\end{evidence}

\begin{claim}[Proposition X.93: Side of third-apotome area is a second apotome of medial]
\label{prop:X.93}
If an area be contained by a rational straight line and a third
apotome, the side of the area is a second apotome of a medial.
\end{claim}
\begin{evidence}[Proof of X.93]
\label{ev:X.93}
Dual of X.56.
\dependson{X.93}{X.56}
\dependson{X.93}{X.87}
\dependson{X.93}{X.92}
\end{evidence}

\begin{claim}[Proposition X.94: Side of fourth-apotome area is a minor]
\label{prop:X.94}
If an area be contained by a rational straight line and a fourth
apotome, the side of the area is a minor.
\end{claim}
\begin{evidence}[Proof of X.94]
\label{ev:X.94}
Dual of X.57.
\dependson{X.94}{X.57}
\dependson{X.94}{X.88}
\dependson{X.94}{X.93}
\end{evidence}

\begin{claim}[Proposition X.95: Side of fifth-apotome area produces rational-plus-medial complement]
\label{prop:X.95}
If an area be contained by a rational straight line and a fifth
apotome, the side of the area is the line producing with a rational
area a medial whole.
\end{claim}
\begin{evidence}[Proof of X.95]
\label{ev:X.95}
Dual of X.58.
\dependson{X.95}{X.58}
\dependson{X.95}{X.89}
\dependson{X.95}{X.94}
\end{evidence}

\begin{claim}[Proposition X.96: Side of sixth-apotome area produces medial-plus-medial complement]
\label{prop:X.96}
If an area be contained by a rational straight line and a sixth
apotome, the side of the area is the line producing with a medial
area a medial whole.
\end{claim}
\begin{evidence}[Proof of X.96]
\label{ev:X.96}
Dual of X.59.
\dependson{X.96}{X.59}
\dependson{X.96}{X.90}
\dependson{X.96}{X.95}
\end{evidence}

\begin{claim}[Proposition X.97: Square on apotome yields first apotome]
\label{prop:X.97}
The square on an apotome straight line applied to a rational
straight line produces as breadth a first apotome.
\end{claim}
\begin{evidence}[Proof of X.97]
\label{ev:X.97}
Inverse of X.91.
\dependson{X.97}{X.85}
\dependson{X.97}{X.91}
\end{evidence}

\begin{claim}[Proposition X.98: Square on first apotome of medial yields second apotome]
\label{prop:X.98}
The square on a first apotome of a medial straight line applied to
a rational straight line produces as breadth a second apotome.
\end{claim}
\begin{evidence}[Proof of X.98]
\label{ev:X.98}
Inverse of X.92.
\dependson{X.98}{X.86}
\dependson{X.98}{X.92}
\dependson{X.98}{X.97}
\end{evidence}

\begin{claim}[Proposition X.99: Square on second apotome of medial yields third apotome]
\label{prop:X.99}
The square on a second apotome of a medial straight line applied to
a rational straight line produces as breadth a third apotome.
\end{claim}
\begin{evidence}[Proof of X.99]
\label{ev:X.99}
Inverse of X.93.
\dependson{X.99}{X.87}
\dependson{X.99}{X.93}
\dependson{X.99}{X.98}
\end{evidence}

\begin{claim}[Proposition X.100: Square on minor yields fourth apotome]
\label{prop:X.100}
The square on a minor applied to a rational straight line produces
as breadth a fourth apotome.
\end{claim}
\begin{evidence}[Proof of X.100]
\label{ev:X.100}
Inverse of X.94.
\dependson{X.100}{X.88}
\dependson{X.100}{X.94}
\dependson{X.100}{X.99}
\end{evidence}

\begin{claim}[Proposition X.101: Square on rational-plus-medial producer yields fifth apotome]
\label{prop:X.101}
The square on the line producing with a rational area a medial whole
applied to a rational straight line produces as breadth a fifth
apotome.
\end{claim}
\begin{evidence}[Proof of X.101]
\label{ev:X.101}
Inverse of X.95.
\dependson{X.101}{X.89}
\dependson{X.101}{X.95}
\dependson{X.101}{X.100}
\end{evidence}

\begin{claim}[Proposition X.102: Square on medial-plus-medial producer yields sixth apotome]
\label{prop:X.102}
The square on the line producing with a medial area a medial whole
applied to a rational straight line produces as breadth a sixth
apotome.
\end{claim}
\begin{evidence}[Proof of X.102]
\label{ev:X.102}
Inverse of X.96.
\dependson{X.102}{X.90}
\dependson{X.102}{X.96}
\dependson{X.102}{X.101}
\end{evidence}

\begin{claim}[Proposition X.103: Commensurable with apotome is apotome]
\label{prop:X.103}
A straight line commensurable in length with an apotome is itself an
apotome and the same in order.
\end{claim}
\begin{evidence}[Proof of X.103]
\label{ev:X.103}
Dual of X.66.
\dependson{X.103}{X.66}
\dependson{X.103}{X.73}
\end{evidence}

\begin{claim}[Proposition X.104: Commensurable with apotome of medial is the same]
\label{prop:X.104}
A straight line commensurable in length with an apotome of a medial
is itself such an apotome of the same order.
\end{claim}
\begin{evidence}[Proof of X.104]
\label{ev:X.104}
Dual of X.67.
\dependson{X.104}{X.67}
\dependson{X.104}{X.74}
\dependson{X.104}{X.103}
\end{evidence}

\begin{claim}[Proposition X.105: Commensurable with minor is minor]
\label{prop:X.105}
A straight line commensurable with a minor is itself a minor.
\end{claim}
\begin{evidence}[Proof of X.105]
\label{ev:X.105}
Dual of X.68.
\dependson{X.105}{X.68}
\dependson{X.105}{X.76}
\dependson{X.105}{X.104}
\end{evidence}

\begin{claim}[Proposition X.106: Commensurable with rational-medial producer is the same]
\label{prop:X.106}
A straight line commensurable with the line producing with a rational
area a medial whole is itself such a line.
\end{claim}
\begin{evidence}[Proof of X.106]
\label{ev:X.106}
Dual of X.69.
\dependson{X.106}{X.69}
\dependson{X.106}{X.77}
\dependson{X.106}{X.105}
\end{evidence}

\begin{claim}[Proposition X.107: Commensurable with medial-medial producer is the same]
\label{prop:X.107}
A straight line commensurable with the line producing with a medial
area a medial whole is itself such a line.
\end{claim}
\begin{evidence}[Proof of X.107]
\label{ev:X.107}
Dual of X.70.
\dependson{X.107}{X.70}
\dependson{X.107}{X.78}
\dependson{X.107}{X.106}
\end{evidence}

\begin{claim}[Proposition X.108: Side of rational minus medial is one of the four irrationals]
\label{prop:X.108}
If from a rational area a medial area be subtracted, the side of the
remaining area arises as one of four irrationals: an apotome, a
first apotome of a medial, a minor, or the line producing with a
rational area a medial whole.
\end{claim}
\begin{evidence}[Proof of X.108]
\label{ev:X.108}
Dual of X.71.
\dependson{X.108}{X.71}
\dependson{X.108}{X.73}
\dependson{X.108}{X.74}
\dependson{X.108}{X.76}
\dependson{X.108}{X.77}
\end{evidence}

\begin{claim}[Proposition X.109: Medial minus rational yields apotome or producer-of-medial]
\label{prop:X.109}
If from a medial area a rational area be subtracted, two other
irrational straight lines arise, namely a first apotome of a medial
or the line producing with a rational area a medial whole.
\end{claim}
\begin{evidence}[Proof of X.109]
\label{ev:X.109}
Variant of X.108.
\dependson{X.109}{X.74}
\dependson{X.109}{X.77}
\dependson{X.109}{X.108}
\end{evidence}

\begin{claim}[Proposition X.110: Medial minus medial yields second apotome of medial or medial-producer]
\label{prop:X.110}
If from a medial area there be subtracted a medial area
incommensurable with the whole, the remaining two irrational straight
lines arise: a second apotome of a medial or the line producing with
a medial area a medial whole.
\end{claim}
\begin{evidence}[Proof of X.110]
\label{ev:X.110}
Variant of X.108 / X.109.
\dependson{X.110}{X.75}
\dependson{X.110}{X.78}
\dependson{X.110}{X.109}
\end{evidence}

\begin{claim}[Proposition X.111: Apotome and binomial are distinct]
\label{prop:X.111}
The apotome is not the same as the binomial.
\end{claim}
\begin{evidence}[Proof of X.111]
\label{ev:X.111}
A binomial has a rational sum of squares plus a medial rectangle; an
apotome has a rational difference of squares minus a medial
rectangle; if they coincided, the two combinations would coincide,
forcing the medial part to be rational --- contradiction.
\dependson{X.111}{X.36}
\dependson{X.111}{X.73}
\end{evidence}

\begin{claim}[Proposition X.112: Square on rational divided by binomial is an apotome]
\label{prop:X.112}
The square on a rational straight line applied to the binomial
straight line produces as breadth an apotome the terms of which are
commensurable with the terms of the binomial and in the same ratio.
\end{claim}
\begin{evidence}[Proof of X.112]
\label{ev:X.112}
$R^2 / (a + b) = a' - b'$ with $a'$, $b'$ in the same ratio as $a$,
$b$.  Verified by direct manipulation of the square-on-binomial
identity.
\dependson{X.112}{X.91}
\dependson{X.112}{X.97}
\end{evidence}

\begin{claim}[Proposition X.113: Square on rational divided by apotome is a binomial]
\label{prop:X.113}
The square on a rational straight line applied to an apotome produces
as breadth a binomial the terms of which are commensurable with the
terms of the apotome and in the same ratio.
\end{claim}
\begin{evidence}[Proof of X.113]
\label{ev:X.113}
Inverse of X.112.
\dependson{X.113}{X.54}
\dependson{X.113}{X.60}
\dependson{X.113}{X.112}
\end{evidence}

\begin{claim}[Proposition X.114: Rectangle on binomial and apotome can be rational]
\label{prop:X.114}
If an area be contained by an apotome and the binomial the terms of
which are commensurable with the terms of the apotome and in the
same ratio, the side of the area is rational.
\end{claim}
\begin{evidence}[Proof of X.114]
\label{ev:X.114}
The rectangle on $(a - b)$ and $(a' + b')$ with $a' = ka$, $b' = kb$
equals $k(a^2 - b^2)$, which is rational.
\dependson{X.114}{X.112}
\dependson{X.114}{X.113}
\end{evidence}

\begin{claim}[Proposition X.115: Medials yield infinitely many irrationals]
\label{prop:X.115}
From a medial straight line there arise irrational straight lines
infinite in number, and none of them is the same with any preceding.
\end{claim}
\begin{evidence}[Proof of X.115]
\label{ev:X.115}
By repeated mean-proportional construction (VI.13) on the medial,
each new line is irrational with respect to all earlier ones (using
the unique-decomposition results X.42--X.47, X.79--X.84).
\dependson{X.115}{VI.13}
\dependson{X.115}{X.21}
\dependson{X.115}{X.114}
\end{evidence}
